\section{预计的困难及解决办法}

\begin{enumerate}
  \item \textbf{GCG优化在非凸损失景观中的收敛不稳定性。} 在智能体场景中，GCG的梯度引导优化面临高度非凸的损失景观，不同随机初始化产生差异极大的攻击效果，甚至随机搜索与梯度优化表现相当。拟通过多种初始化策略（随机初始化与语义模板初始化相结合）、增大批量搜索规模、增加优化迭代次数来缓解收敛不稳定问题，并以多次独立运行的统计结果（如Security-at-N指标）而非单次最优结果来报告攻击性能。

  \item \textbf{TAP中攻击者LLM拒绝行为对搜索的干扰。} 攻击者LLM在生成注入候选时可能因安全对齐而拒绝生成攻击内容，且拒绝行为在树搜索的分支中具有传播，一旦某分支出现拒绝，后续迭代倾向于持续拒绝，导致搜索空间坍缩。拟通过精心设计系统提示（明确安全研究上下文、提供角色扮演框架）、设计多样化的变异策略模板来降低拒绝率，同时在剪枝阶段优先保留无拒绝分支，并在分支多样性不足时引入重新初始化机制。

  \item \textbf{LLM-as-judge评估准确性不足导致优化偏差。} 评估器LLM对攻击成功与否的判断准确性因目标模型而异（如Qwen上仅约50\%准确率），导致假阴性剪枝丢弃有希望的攻击候选，假阳性评分浪费优化预算。拟针对不同目标模型调整评估器阈值（对低准确率模型采用更保守的剪枝策略）、增加每候选的可靠性测试评估次数、并结合AgentDojo的确定性检查函数进行后验证校准。

  \item \textbf{跨模型攻击迁移率极低。} 在开源小模型（Qwen3-4B、Gemma-3-4B）上优化的攻击几乎完全无法迁移到先进闭源模型（GPT-5），成功率低于2\%。这一根本性挑战源于模型架构、训练数据和安全对齐策略的差异。拟从实验分析角度系统量化迁移率差距，探索通用攻击中语义结构注入（如TAP生成的自然语言注入）相对于token级优化注入的迁移优势，并通过消融实验识别影响迁移性的关键因素。

  \item \textbf{实验规模与计算资源约束。} 80个任务对 $\times$ 多种攻击方法 $\times$ 多种配置 $\times$ 多次独立运行的实验矩阵规模庞大，GCG优化需要大量GPU梯度计算，TAP优化需要大量API调用。拟通过合理的实验子集选择（基于任务难度和长度进行定性筛选代表完整基准的子集）、GCG批量前向传播优化、TAP并行评估加速等策略控制计算开销，确保实验在可行时间内完成。
\end{enumerate}
